% -----------------------------------------------------------------------------
% Conteudo AODV-EN. Template base: Prof. Wesley Pacheco Calixto (IFG).
% -----------------------------------------------------------------------------
\addcontentsline{toc}{section}{Resumo}

\section*{Resumo}

As redes mesh sem fio de baixo custo construídas sobre o protocolo ESP-NOW carecem de
suporte nativo a roteamento multi-hop, o que limita seu alcance e escalabilidade. Este
trabalho propõe, implementa e avalia o AODV-EN, uma adaptação do protocolo Ad-hoc
On-Demand Distance Vector (AODV, RFC 3561) para operação sobre ESP-NOW em módulos ESP32,
contemplando endereçamento por MAC, disseminação de requisições de rota por
\emph{flooding} controlado sobre \emph{broadcast} de enlace, manutenção de rotas com
números de sequência e confirmação fim-a-fim. Adotando o paradigma \emph{Design Science
Research}, o núcleo do protocolo foi implementado em linguagem C como componente
reutilizável e comparado a um algoritmo de referência baseado em \emph{flooding},
medindo-se \emph{Packet Delivery Ratio} (PDR), latência fim-a-fim, \emph{Normalized
Routing Load} (NRL) e consumo energético estimado. No cenário avaliado em hardware (três
ESP32, seis repetições), os dois algoritmos apresentaram confiabilidade equivalente
(PDR de 98,93\% para o AODV-EN e 98,77\% para o \emph{flooding}), com vantagem do AODV-EN
em latência (60,0 ms estáveis contra 62 a 100 ms) e em consumo energético estimado
(aproximadamente 21\% inferior), ao custo de uma carga de controle normalizada de 0,77.
A análise por simulação evidenciou que a desvantagem de ocupação de canal do
\emph{flooding} cresce com a escala da rede, com cruzamento de eficiência em torno de
9 a 11 nós. Conclui-se que o roteamento reativo do AODV-EN é vantajoso à medida que a
rede cresce em tamanho e diâmetro, viabilizando redes mesh padronizadas, de baixo custo
e eficientes em energia sobre ESP-NOW.

\vspace{1cm}

\noindent
\textbf{Palavras-chave:} redes mesh; ESP-NOW; ESP32; AODV; roteamento ad-hoc;
Internet das Coisas.
